% !TEX root = ../../tree_notes.tex
% =====================================================================
%  Chapter 5 opener -- "Chapter at a Glance", reading paths + ladder
%  staircase, the running example G (reused from Ch.4) now stamped with
%  tin/tout, roadmap, topic table, recognition guide, constraint table,
%  notation.  NO numbered boxes here (they would be numbered 5.0.x).
% =====================================================================

\begin{chaptertheorybox}{Chapter at a Glance}
Chapters~3 and~4 taught you to \emph{walk} a tree and to \emph{store} it. This chapter
makes the leap that turns a tree into something a \emph{segment tree can eat}: one DFS
hands every node two integers, a \term{birth time} $\tin(v)$ and a \term{death time}
$\tout(v)$, and suddenly an entire \emph{subtree} -- a branching, recursive shape -- is
just a \term{contiguous slice} $[\tin(v), \tout(v))$ of a flat array. That single idea,
the \term{Euler tour} (or \term{flattening}), is the gateway from "tree" to "array with
range queries", and it is the reason the Fenwick, segment-tree and heavy-light chapters
exist. The three decisions this chapter forces on you:

\begin{itemize}[nosep]
  \item \textbf{When did I enter and leave each node?} A DFS that stamps a counter on
        entry (and again on exit) produces $\tin/\tout$; the containment test
        "$\tin(u) \le \tin(v) < \tout(u)$" answers \emph{is $u$ an ancestor of $v$?} in
        $O(1)$, and the same times spell the tree as a \term{balanced bracket sequence}.
        \hfill \lv{2}--\lv{3}
  \item \textbf{How do I put a subtree on a line?} Order the array by $\tin$; then
        $\subtree(v)$ is exactly the block $[\tin(v), \tin(v){+}\size[v])$. Any subtree
        sum / max / count becomes a \emph{range query} and any "paint the whole subtree"
        becomes a \emph{range update}.   \hfill \lv{3}--\lv{4}
  \item \textbf{What about paths and distances?} The same tour, read differently, gives
        the root-to-node \emph{path as a stack}, the \emph{depth sequence} for level
        problems, and the $2n{-}1$ walk whose first-occurrence array turns
        $\dist(u,v)$ and LCA into a range-minimum query.   \hfill \lv{3}--\lv{5}
\end{itemize}

\textbf{You will learn to}
\begin{itemize}[nosep]
  \item stamp $\tin/\tout$ with one DFS and test ancestry in $O(1)$;          \hfill\lv{2}--\lv{3}
  \item recognise the bracket / contour view of a tree and read it back;       \hfill\lv{3}
  \item flatten a tree to an array so a subtree is a range, then answer or update
        that range (the hand-off to the query chapters);                       \hfill\lv{3}--\lv{4}
  \item aggregate along root-to-node paths with an in/out sweep and per-level arrays;
        \hfill\lv{3}--\lv{4}
  \item build the $2n{-}1$ Euler tour, index it by first occurrence, and reduce
        distance / LCA to a range-minimum query (proved in the LCA chapter).   \hfill\lv{4}--\lv{5}
\end{itemize}

\textbf{Prerequisites.} The four DFS orders of Chapter~3 (\chref{c03}) and the
parent-guarded rooting DFS of Chapter~4 (\secref{c04-rooting}, \coderef{c04-root}) -- this
chapter \emph{stamps times during} that DFS; it does not re-derive traversal or rooting.
Subtree aggregation is \secref{c02-bottomup}.

\textbf{In contests and interviews.} Whenever a problem says "for each node, query/modify
its whole subtree" or "the distance between two nodes", the intended solution almost always
starts \emph{here}: run one DFS, record $\tin/\tout$, and reduce the tree question to an
array question you already know how to solve. This chapter is that reduction, end to end.
\end{chaptertheorybox}

% ---------------- Station 1b: reading paths + ladder position ----------------
\begin{chaptertheorybox}{How to Read This Chapter (two paths up the ladder)}
Level badges $\lv1$--$\lv6$ as before: $\lv2$ is \emph{``stamp times, test ancestry''},
$\lv3$ is the solid interview medium, $\lv5$ is LC Hard with a tour-plus-RMQ or a
flattened-split on top.

\textbf{Beginner path ($\lv2$--$\lv3$).} Read \secref{c24-times} fully -- it fixes
$\tin/\tout$, the $O(1)$ ancestor test and the bracket view. Then \secref{c24-flatten} up
through "subtree $=$ range". Take every $\lv2$--$\lv3$ solved problem. You may \emph{skip
on a first pass}: the depth-sequence/path material of \secref{c24-path}, the $2n{-}1$ tour
and RMQ preview of \secref{c24-dist}, every $\lv4{+}$ variation, the hardening notes and the
Summit. At the end of this path you can flatten any tree and answer "is this an ancestor /
what is this subtree's total" instantly.

\textbf{Full path.} Everything, in order -- it ends with the scaled summit (LC Hard 2322
\emph{Minimum Score After Removals on a Tree}: root it, take $\tin/\tout$, and read every
edge-deletion as splitting one interval into two) and the reduction gallery showing that
ancestry, subtree-range and distance are the \emph{same tour} seen from three angles. The
genuine $\lv5$--$\lv6$ query machinery -- Fenwick/segment trees over the flat array, HLD's
path decomposition, the RMQ-as-LCA proof -- lives in the query and LCA chapters
(\chref{c14}, \chref{c13}, \chref{c16}, \chref{c11}).
\end{chaptertheorybox}

\begin{figure}[H]
\centering
\begin{tikzpicture}[diagram, x=9mm, y=7mm]
  \foreach \st/\lab in {1/lv1,2/lv2,3/lv3,4/lv4,5/lv5,6/lv6} {
    \fill[\lab!22, draw=\lab, line width=0.8pt] (\st-1,0) rectangle (\st,\st*0.52);
    \node[font=\scriptsize\bfseries, text=\lab] at (\st-0.5,\st*0.52-0.42) {L\st};}
  \node[annot, anchor=south west, font=\scriptsize]
    at (0.05,0.62) {5.1 stamp times, ancestry};
  \node[annot, anchor=south west, font=\scriptsize]
    at (1.05,1.14) {5.1 bracket / contour};
  \node[annot, anchor=south west, font=\scriptsize]
    at (2.05,1.66) {5.2 subtree-as-range};
  \node[annot, anchor=south west, font=\scriptsize]
    at (3.05,2.18) {5.2 hand-off to range structures};
  \node[annot, anchor=south west, font=\scriptsize, text=lv5]
    at (4.05,2.70) {5.3 paths / depth seq, 5.4 2n-1 tour};
  \node[annot, anchor=south west, font=\scriptsize, text=lv5]
    at (5.05,3.22) {scaled summit (2322)};
\end{tikzpicture}
\caption*{Chapter~5 on the six-step ladder: the two flattening topics carry you from raw
timestamps to "subtree $=$ range" ($\lv4$); the path/distance topics and the scaled summit
(2322) reach $\lv5$. Beginner path $=$ steps 2--3.}
\end{figure}

% ---------------- The running example ----------------
\subsection*{The Running Example}

This chapter has \emph{one} running example, and it is Chapter~4's general contest tree
$G$ -- the very tree whose $\parent$, $\depth$ and $\size$ arrays you already know
(\figref{c04-running-g}). Reusing it is the point: the same rooted tree, seen through
timestamps, becomes an array. Every figure, worked example and trace here uses $G$ and
nothing else. Rooted at $0$, visited children in ascending order, the DFS walk is
$0, 1, 4, 5, 2, 3, 6$.

\begin{figure}[H]
\centering
\begin{tikzpicture}[diagram]
  \node[troot] (n0) at (0,0) {$0$};
  \node[tnode] (n1) at (-3,-1.4) {$1$};
  \node[tnode] (n2) at (0,-1.4) {$2$};
  \node[tnode] (n3) at (3,-1.4) {$3$};
  \node[tleaf] (n4) at (-3.9,-2.8) {$4$};
  \node[tleaf] (n5) at (-2.1,-2.8) {$5$};
  \node[tleaf] (n6) at (3,-2.8) {$6$};
  \draw[tedgehl,->] (n0) -- (n1); \draw[tedgehl,->] (n0) -- (n2); \draw[tedgehl,->] (n0) -- (n3);
  \draw[tedgehl,->] (n1) -- (n4); \draw[tedgehl,->] (n1) -- (n5); \draw[tedgehl,->] (n3) -- (n6);
  \node[annot, align=left, font=\scriptsize] at (-6.3,-0.3) {$\tin$/$\tout$ (this chapter)\\
    $\tin=[0,1,4,5,2,3,6]$\\$\tout=[7,4,5,7,3,4,7]$};
  \node[annot, align=left, font=\scriptsize] at (4.4,-0.3) {carried from Ch.4\\$\size=[7,3,1,2,1,1,1]$\\$\depth=[0,1,1,1,2,2,2]$};
\end{tikzpicture}
\caption{Running example $G$ (rooted at $0$). The DFS preorder is
$0,1,4,5,2,3,6$; numbering nodes \emph{by entry time} gives $\tin$ and, since
$\tout(v)=\tin(v)+\size[v]$, $\tout$ as above. Node $1$'s subtree is the block
$[\tin(1),\tout(1))=[1,4)$ of the preorder array $[\,0,1,4,5,2,3,6\,]$, i.e.\ the
nodes $\{1,4,5\}$ -- the whole chapter in one picture.}\label{fig:c24-running}
\end{figure}

\medskip
\noindent Keep these four lines for $G$ in mind -- they are re-derived, never asserted, in
\secref{c24-times}:
\begin{center}
\renewcommand{\arraystretch}{1.2}
\begin{tabular}{@{}p{0.30\linewidth} p{0.62\linewidth}@{}}
\toprule
\textbf{View of $G$} & \textbf{Value (rooted at $0$, children ascending)} \\
\midrule
preorder (visit order)          & $0,\,1,\,4,\,5,\,2,\,3,\,6$ \\
$\tin[i]$ (entry time $=$ preorder rank) & $[0,\,1,\,4,\,5,\,2,\,3,\,6]$ \\
$\size[i]$ (subtree node count) & $[7,\,3,\,1,\,2,\,1,\,1,\,1]$ \\
$\tout[i]=\tin[i]+\size[i]$ (exit) & $[7,\,4,\,5,\,7,\,3,\,4,\,7]$ \\
\bottomrule
\end{tabular}
\end{center}

% ---------------- Roadmap diagram ----------------
\begin{figure}[H]
\centering
\begin{tikzpicture}[diagram, node distance=5mm,
    top/.style={boxnode, fill=primaryblue!10, draw=primaryblue, text width=25mm,
                minimum height=11mm, font=\small},
    res/.style={annot, text width=25mm, anchor=north}]
  \node[top] (t1) {\textbf{5.1} Times \& Bracket};
  \node[top, right=of t1] (t2) {\textbf{5.2} Subtree $=$ Range};
  \node[top, right=of t2] (t3) {\textbf{5.3} Path \& Depth};
  \node[top, right=of t3] (t4) {\textbf{5.4} Distance \& RMQ};
  \foreach \a/\b in {t1/t2, t2/t3, t3/t4} \draw[flow] (\a) -- (\b);
  \node[res] at ([yshift=-2mm]t1.south) {tin/tout; O(1) ancestry; balanced brackets; contour};
  \node[res] at ([yshift=-2mm]t2.south) {flatten to array; range sum/max/count; hand-off to BIT/seg};
  \node[res] at ([yshift=-2mm]t3.south) {root-path as stack; level/depth arrays; in/out sweep};
  \node[res] at ([yshift=-2mm]t4.south) {2n-1 tour; first-occurrence; dist formula; RMQ-LCA preview};
\end{tikzpicture}
\caption*{The four topics, left to right, all from \emph{one DFS}. 5.1 asks \emph{when do I
enter/leave}; 5.2 asks \emph{what does a subtree look like on a line}; 5.3 asks \emph{what
does a path look like}; 5.4 asks \emph{what does a distance look like}. Each is the same
tour read one way.}
\end{figure}

% ---------------- Topic table ----------------
\begin{center}
\renewcommand{\arraystretch}{1.25}
\setlength{\tabcolsep}{4pt}
\begin{tabular}{@{}p{0.24\linewidth} p{0.36\linewidth} c p{0.23\linewidth}@{}}
\toprule
\textbf{Topic} & \textbf{Key results} & \textbf{Levels} & \textbf{Solved problems} \\
\midrule
\secref{c24-times}\ \ Times \& Bracket
  & $\tin/\tout$ from one DFS; $O(1)$ ancestor test; bracket/contour encoding & $\lv2$--$\lv3$
  & \probref{c24-throne} \\
\addlinespace
\secref{c24-flatten}\ \ Subtree as Interval
  & flatten by $\tin$; subtree $=$ range; range sum/max/count; hand-off to query trees & $\lv3$--$\lv4$
  & \probref{c24-fuel} \\
\addlinespace
\secref{c24-path}\ \ Path \& Depth Sequence
  & root-path as stack; level/depth arrays; per-level aggregation & $\lv2$--$\lv3$
  & \probref{c24-levelsum}, \probref{c24-kthlevel} \\
\addlinespace
\secref{c24-dist}\ \ Distance \& RMQ Preview
  & $2n{-}1$ tour; first-occurrence array; $\dist(u,v)$; RMQ-as-LCA teaser & $\lv4$--$\lv5$
  & \probref{c24-pathdir}, \probref{c24-leafpairs} \\
\bottomrule
\end{tabular}
\end{center}

% ---------------- Problem recognition guide ----------------
\subsection*{Problem Recognition Guide --- When You See \emph{this}, Think \emph{that}}

Flattening problems announce themselves by asking about a \emph{whole subtree} or a
\emph{path/distance} many times. Match the phrasing; every row points at the section that
builds the tool.

\begin{center}
\renewcommand{\arraystretch}{1.3}
\setlength{\tabcolsep}{4pt}
\begin{tabular}{@{}p{0.42\linewidth} p{0.35\linewidth} l@{}}
\toprule
\textbf{When the problem says\ldots} & \textbf{\textcolor{deepblue}{\textbf{Think\ldots}}} & \textbf{Section} \\
\midrule
``is $u$ an ancestor of $v$?'' / ``the $k$-th successor of a node'' / a subtree ordering
  & stamp $\tin/\tout$; answer ancestry in $O(1)$ & \secref{c24-times} \\
\addlinespace
``a string of nested parentheses'' / ``reconstruct from brackets'' / ``$\le$ depth''
  & the bracket / contour view of the same tour & \secref{c24-times} \\
\addlinespace
``for each node, \emph{sum/count/max over its whole subtree}'' (many queries)
  & flatten by $\tin$; subtree $=$ range $[\tin,\tin{+}\size)$; range structure & \secref{c24-flatten} \\
\addlinespace
``add $x$ to every node in the subtree of $v$'', then read a node
  & range \emph{update} on the flat array (defer the Fenwick/seg chapter) & \secref{c24-flatten} \\
\addlinespace
``the sum of values on the \emph{path from the root} to $v$'' / ``per-level answer''
  & carry a path stack / level array down the DFS & \secref{c24-path} \\
\addlinespace
``the \emph{distance} between nodes $u$ and $v$'' / ``path between two nodes''
  & $\dist(u,v)=\depth u+\depth v-2\depth{\lca}$; LCA via the $2n{-}1$ tour & \secref{c24-dist} \\
\bottomrule
\end{tabular}
\end{center}

\begin{center}
\renewcommand{\arraystretch}{1.25}
\begin{tabular}{@{}p{0.27\linewidth} p{0.45\linewidth} p{0.20\linewidth}@{}}
\toprule
\textbf{Constraint you see} & \textbf{What it forces} & \textbf{Chapter-5 tool} \\
\midrule
many ``over the subtree of $v$'' queries & do \emph{not} DFS per query -- flatten once, answer by range & \secref{c24-flatten} \\
$N$ up to $2\cdot10^{5}$, point/subtree updates & flatten $+$ a Fenwick/segment tree (\chref{c14}, \chref{c13}) & \secref{c24-flatten} \\
"distance between two nodes" asked fast & LCA; here the tour + RMQ \emph{preview}, proved \chref{c11} & \secref{c24-dist} \\
tree may be a chain (depth $\approx n$) & recursion risks RE -- do the tour \emph{iteratively} & \secref{c24-times} \\
\bottomrule
\end{tabular}
\end{center}

\begin{notebox}[Not this chapter?]
If you merely need to \emph{visit} nodes in an order, that is Chapter~3 (\chref{c03}); this
chapter uses the order but adds timestamps. If the tree is not yet \emph{rooted} (raw
undirected edges), root it first in Chapter~4 (\secref{c04-rooting}). If the actual task is
to \emph{answer range queries with updates}, the range structure itself is Fenwick
(\chref{c14}) or segment (\chref{c13}) -- we only reduce the tree to a range and hand it
over. If you need LCA \emph{proved} and used for path queries, that is \chref{c11} and
\chref{c16}. If you need to treat subtrees \emph{offline in blocks}, that is Mo on trees
(\chref{c21}); \emph{per-node subtree merging} is DSU-on-tree/sack (\chref{c28}); \emph{subtree
identity} is hashing (\chref{c22}). This chapter owns \emph{turning the tree into a line}.
\end{notebox}

% ---------------- Notation ----------------
\subsection*{Notation Used in This Chapter}

\begin{center}
\renewcommand{\arraystretch}{1.2}
\begin{tabular}{@{}p{0.20\linewidth} p{0.27\linewidth}@{\quad}p{0.17\linewidth} p{0.27\linewidth}@{}}
\toprule
\textbf{Symbol} & \textbf{Meaning} & \textbf{Symbol} & \textbf{Meaning} \\
\midrule
$G$               & the running general tree (\figref{c24-running}) & $\tin(v)$ & DFS \emph{entry} time of $v$ (preorder rank) \\
$\tout(v)$        & $\tin(v)+\size[v]$, the \emph{exit} time       & $\size[v]$ & node count of $\subtree(v)$ (from Ch.4) \\
$\depth(v)$       & edges from the root to $v$                    & $\subtree(v)$ & the tree rooted at $v$ \\
$\flat[\,]$       & the flattened array, indexed by $\tin$        & $\dist(u,v)$ & number of edges on the $u$--$v$ path \\
$\lca(u,v)$       & lowest common ancestor (proved \chref{c11})   & $\mathcal{T}$ & the $2n{-}1$ Euler tour sequence \\
$\mathrm{first}[\,]$ & first-occurrence index in $\mathcal{T}$    & $\lv{3}$ & difficulty badge (runs $\lv{1}$--$\lv{6}$) \\
\bottomrule
\end{tabular}
\end{center}

As in Chapters~1--4, $\depth$ and $\dist$ are counted in \textbf{edges}, general trees are
stored as an adjacency list (\code{adj}, a \code{vector} of \code{vector<int>}) with a parent
guard, every listing is complete $\cpp$ (the standard header and the \code{TreeNode} where the tree is
binary), and node numbering follows the problem: LeetCode binary trees are \emph{unindexed}
(pure pointers), contest general trees are \emph{$0$- or $1$-indexed} by node id. Times are
$0$-indexed here; a $1$-indexed clock works identically.
